\chapter{Introduction}


Adversarial planning against non-cooperative agents has become a relevant area of research in real-world autonomous systems and robotics problems. Adversarial problems arise when an autonomous agent must make decisions while accounting for the actions of another intelligent and non-cooperative opponent with opposing objectives. These sorts of problems appear in various multi-agent planning settings, such as pursuit-evasion, target-attacker-defender (\acs{tad}) asset defense, and robotic guidance tasks. Examples of these sorts of games are provided in the literature review in chapter \ref{chap:Lit_Review}. In any of these, a successful outcome for an agent depends not only on their actions and the system dynamics, but also on the evolving strategy of its opponents. Traditional optimal control frameworks to solve these kinds of problems tend to assume certain information strategies which can be jeopardized against an agent that's following a shifting, actively adversarial strategy coming from a versatile, intelligent opponent. 

In these cases, classical control architectures may prove to be brittle against these highly dynamic strategies. Namely, this can happen when the opponent doesn't follow an assumed strategy, information is limited, or if solving for a pursuit-evasion game through online optimization methods becomes computationally demanding. When coupled with estimation and safety constraints in the loop, these  difficulties can become even more pronounced.

Inspired by this dilemma, we chose to develop a reinforcement-learning (\acs{rl}) based approach for planning against these sorts of adversarial opponents. Instead of solving an optimization game online on every timestep, the \acs{rl} framework concentrates most of its online computational burden on training, therefore producing a deterministic policy with fast inference speeds during rollout. The benefit behind these \acs{rl}-based architectures is that, when properly trained, these can produce a single, deterministic policy that can actively react and adapt to a wide range of scenarios it's been trained against.

One of the arguments behind why research and industry have not widely adopted \acs{rl} algorithms within their autonomy/\acs{gnc} work comes down to their lack of formal guarantees for stability and convergence. This lack of guarantees can be attributed to the fact that these \acs{rl}-based algorithms try to optimize against highly non-convex objectives while using gradients from noisy trajectories, approximate the value functions they learn during training, are prone to overfit to the data distribution that they're trained on, and struggle to interpolate to data outside this distribution. In other words, these factors can lead to a noisy stochastic gradient ascent/descent process that finds locally optimal solutions instead of the globally optimal solution that may not be guaranteed to work in regimes outside the training dataset. In fact, the development of better \acs{rl} algorithms that can account for these mishaps present in the current state-of-the-art \acs{rl} is a core area of research in world of theoretical deep \acs{rl} in this day and age. 

Regardless, algorithms such as Proximal Policy Optimization (\acs{ppo}, \cite{schulman2017ppo}), an on-policy algorithm created by OpenAI, remain one of the most widely used algorithms in \acs{rl} research for engineering applications (\cite{francois2018introduction}). In simple terms, \acs{ppo} is based on the principles of stochastic gradient ascent on a clipped surrogate objective function to be optimized. The benefits of these \acs{ppo}-based architectures are that they're relatively simple to set up, they reuse each rollout batch on multiple gradient steps, are easy to parallelize, and, most importantly, provide much more stability than vanilla policy gradient methods. For-controls based engineering applications, which is what we'll be studying in this work, perhaps the biggest benefit of \acs{ppo} is that it can be used for both continuous and discrete action spaces.



With all of the above in mind, we set off to study a reinforcement-learning framework to train, model, and rollout robust adversarial policies for orbital zero-sum games in Low Earth Orbit (\acs{leo}). In order to better understand how these architectures could interact with classical autonomy methods, we were most interested in studying how this framework could learn while subject to a closed-loop aerospace guidance, navigation, and control architecture (\acs{gnc}). For our own applications, we wanted to solve a few questions in this research:

\begin{itemize}
    \item Can an \acs{rl} architecture using \acs{ppo} learn coherent policies for continuous control of orbital spacecraft in \acs{leo} in zero-sum adversarial scenarios?
    \item Can an \acs{rl} architecture learn coherent policies while interacting subject to the constraints of a classical controller in the loop?
    \item Can an \acs{rl} architecture learn coherent policies if, instead of having access to the opponent's full state, it had partial observability and had to rely on measurements in the loop?
    \item Given boundaries imposed by the controls and estimation architectures, can an \acs{rl} architecture maintain coherent performance on observations outside of the training dataset even if it was never trained against some of that data?
    \item Can an \acs{rl} architecture do all of the above at once, and still provide convergent and high-performing policies despite the stochasticity?
    \item Can an \acs{rl} architecture trained with \acs{hcw} linear time-invariant orbital relative dynamics transfer to elliptical linear time-varying dynamics that were not seen during training?
\end{itemize}


With this in mind, we propose the Phased Adversarial Reinforcement Learning and Behavior Cloning (\acs{parlbc}) method for training and rollout of \acs{rl}-enabled Guidance, Navigation, and Control (\acs{gnc}) aerospace applications. The overarching goal of the \acs{parlbc} framework, as a whole, is  to develop a phased \acs{ppo}-based learning framework that makes it possible to apply \acs{rl} formulations in a dynamic adversarial problem with physical constraints. Our method uses dense step per geometry-based rewards and phased opponent freezing to mitigate non-stationarity encountered in multi-agent training and aid learning stability. \acs{parlbc} learns while subject to the constraints of a classical velocity-saturating \acs{cbfqp} controller and, optionally, partial observability of the opponent's state via a bearings-only Extended Kalman Filter (\acs{ekf}), which is achieved by using behavior cloning.


Our Monte Carlo evaluation results show that \acs{parlbc} produces highly performing policies in a constrained, zero-sum orbital asset defense game with equal spacecraft characteristics. The defender policies manage an aggregate success rate of up to 79.9\% across four learning phases in the fully observable setting, where the opponent is also guided by an \acs{rl}-trained policy. For the Extended Kalman Filter-based policies, where neither agent has direct access to the opponent's state and instead has to estimate it through a bearings-only \acs{ekf}, the defender succeeds in up to 74.7\% of the cases, a 6.51\% relative drop in performance relative to the fully observable policies. We performed all our evaluations across both \acs{hcw} and elliptic dynamics, even though the policies were never exposed to the latter during training, from which we found at most a 1.29\% relative change between results across dynamics models. Furthermore, these results were achieved while providing fast rollout and compute inference times of 0.145 ms median per timestep on the full-state policies and 3.91 ms per timestep on the \acs{ekf}-based policies.

Our findings open the way for the development of highly computationally efficient \acs{rl}-based methods for autonomous systems research and application, and a discussion of potential directions for future work is provided.

For guidance on the terminology used throughout this writing, please refer to Appendix \ref{appendix:glossary} which serves as a glossary of some of the key terms used throughout this research.
